\section*{Chapitre \ref{ch:hf:adverse-selection-and-toxicity} --- Sélection adverse et toxicité}

\begin{solution}{exo:hf:adverse-selection-and-toxicity:1}
$0.32\times(-0.95)+0.68\times0.55=-0.304+0.374=+0.07$ ticks par action.
\end{solution}

\begin{solution}{exo:hf:adverse-selection-and-toxicity:2}
Prévision $1.432-0.029\times2-1.795\times0.8-0.072\ln1.5+0=-0.09$~: le score vaut 0.09. Le teneur de marché s'efface au seuil de 0 et maintient le cours acheteur à 0.2.
\end{solution}

\begin{solution}{exo:hf:adverse-selection-and-toxicity:3}
$0.2\times10\,000\times\$0.01=\$20$.
\end{solution}

\begin{solution}{exo:hf:adverse-selection-and-toxicity:4}
Chaque changement d'état d'un côté annule un ordre et chaque retour en envoie un nouveau~; le teneur de marché de base ne recote que lorsque le meilleur prix bouge.
\end{solution}

\begin{solution}{exo:hf:adverse-selection-and-toxicity:5}
Dans un carnet à un tick, le prochain mouvement de prix survient lorsqu'une file au meilleur cours est vidée~; un côté qui ne détient qu'une faible part des deux files est celui qui va être
vidé, ce qui se produit quand des traders informés le prennent. Les autres variables sont des versions plus faibles de la même information.
\end{solution}

\begin{solution}{exo:hf:adverse-selection-and-toxicity:6}
$+0.83$ ticks par action (erreur type 0.25), contre $-0.86$ (0.12) pour le pire cinquième~: un écart de 1.69 tick, plusieurs fois les erreurs types.
\end{solution}

\begin{solution}{exo:hf:adverse-selection-and-toxicity:7}
Avec hystérésis~: 716 messages par séance contre 952, 9\,767 actions contre 10\,733, une part d'informés de 21.2~\% contre 17.9~\%, un mark-out à dix secondes
de 0.102. Elle réduit d'un quart le nombre de messages, mais ne baisse pas la toxicité~: un côté maintenu à l'écart plus longtemps manque aussi les bons moments et revient
en queue de sa file.
\end{solution}

\begin{solution}{exo:hf:adverse-selection-and-toxicity:8}
Le contenu prédictif du VPIN provient surtout de l'intensité de négociation (Andersen et Bondarenko), et une mesure de tension à l'échelle du marché ne dit rien sur
lesquelles de ses propres cotations vont subir une sélection adverse. Élargir partout revient à abandonner les bons flux avec les mauvais~; il faut attribuer un score à ses propres exécutions, côté par côté, et agir sur les
côtés que le score signale.
\end{solution}

\begin{solution}{pb:hf:adverse-selection-and-toxicity:1}
\begin{enumerate}
\item Bourse anonyme~: plateforme, type d'ordre, taille, heure de la journée, état du carnet~; demande de cotation (RFQ)~: la contrepartie elle-même.
\item Les ordres à cours limité acheminés vers des plateformes versant de fortes remises ont été exécutés moins bien selon plusieurs mesures de qualité, la sélection adverse parmi elles.
\item Informés $-0.95$ (erreur type 0.14), non informés $+0.55$ (0.10) ticks par action à dix secondes~; 32~\% des exécutions sont informées.
\item Les exécutions informées représentent un tiers des exécutions et perdent presque deux fois ce que les autres rapportent~: $+0.07$ au total.
\item Une estimation de la toxicité du flux à partir du déséquilibre des ordres en temps de volume~; un mauvais prédicteur de la volatilité, maximal seulement après le flash crash, son contenu
est mécanique, porté par l'intensité de négociation.
\item Voir la \cref{def:hf:adverse-selection-and-toxicity:score}~; flux contraire au côté, part de la file du côté opposé, temps calme, une fourchette large.
\item $1.43-0.029\,\text{flux}-1.79\,\text{part}-0.072\ln(1+\text{calme})+0.028\,\text{large}$~; 0.19 dans l'échantillon (in-sample), 0.18 hors échantillon (out-of-sample).
\item Les mark-outs réalisés croissent d'un cinquième à l'autre, de $-0.86$ à $+0.83$~: le score classe correctement les exécutions hors échantillon (out-of-sample).
\item Voir la \cref{def:hf:adverse-selection-and-toxicity:widening}~; l'évanouissement de la cotation retire celle-ci, le skew déplace les deux cotations, refuser les prix ou rejeter une
contrepartie connue.
\item Évanouissement et élargissement d'un tick.
\item Aucun~: 19\,167 actions, 29.4~\%, $-0.005$~; évanouissement~: 6\,550, 18.3~\%, 0.136~; élargissement~: 6\,683, 20.0~\%, 0.112.
\item Les erreurs types des mark-outs (de 0.10 à 0.22) sont aussi grandes que les écarts.
\item Part d'informés de 29.4~\% à 17.9~\%, mark-out à dix secondes de $-0.005$ à 0.214 ticks par action (erreur type 0.146), volume de 19\,167 à
10\,733 actions par séance~: 44~\% sont abandonnés.
\item Près de zéro, le score franchit le seuil dans les deux sens~; chaque franchissement envoie l'ordre en queue de sa file, là où les exécutions sont les pires.
\item Le teneur de marché s'efface trop rarement~: à 0.4, la part d'informés est de 24.1~\%, à 0.6 de 30.1~\%.
\item Moins d'exécutions informées pendant la fenêtre d'annonce (23.6~\% contre 34.7~\% auparavant)~: le flux de bruit augmente avec les annonces~; la toxicité varie dans le temps.
\item Recoter à petite taille ou à un prix normal pendant une période, calculer les mark-outs de ses exécutions, et rétablir la cotation si elles ne sont pas moins bonnes que la moyenne.
\item Le choix de la profondeur redevient central~: l'élargissement devient une réponse continue (la moitié du mouvement attendu du chapitre~4), et l'évanouissement est rarement nécessaire.
\item Tarifer la contrepartie (une cotation plus large ou un taux de réussite plus faible), ou rejeter sous last look.
\item Pour décider, côté par côté et instant par instant, où afficher la liquidité.
\end{enumerate}
\end{solution}

\begin{solution}{iq:hf:adverse-selection-and-toxicity:1}
Contreparties différentes (le flux informé se concentre sur la plateforme)~; frais différents (les plateformes à remises attirent les preneurs toxiques, comme l'ont établi Battalio, Corwin et
Jennings)~; latence différente (on est plus lent sur cette plateforme et on subit la sélection adverse). Découper les mark-outs par type d'ordre et par moment, comparer les délais d'exécution à sa propre
latence, et vérifier si l'écart évolue avec la grille tarifaire.
\iqlookfor{plusieurs hypothèses et un moyen de les séparer.}
\end{solution}

\begin{solution}{iq:hf:adverse-selection-and-toxicity:2}
Biais de sélection~: les exécutions sont plus fréquentes lorsque les cotations sont périmées ou lorsque le score était bas, et les cotations retirées n'ont pas de mark-out. Randomiser une partie de la
politique, pondérer par la probabilité d'être exécuté, ou modéliser conjointement l'exécution et le mark-out.
\iqlookfor{reconnaître le biais de sélection et un remède.}
\end{solution}

\begin{solution}{iq:hf:adverse-selection-and-toxicity:3}
Ça peut être le cas~: les mark-outs sont très bruités, donc une faible corrélation trie encore les exécutions d'un tick ou plus. Ce qui compte, c'est l'écart des mark-outs réalisés
entre les classes de score hors échantillon (out-of-sample), le volume abandonné, et le coût de l'action (messages, place perdue dans la file).
\iqlookfor{raisonner en classes et en coûts, et non en $R^2$.}
\end{solution}

\begin{solution}{iq:hf:adverse-selection-and-toxicity:4}
L'ordre est annulé et ressaisi à chaque franchissement, ce qui fait perdre sa place dans la file et sature la passerelle. Ajouter une hystérésis ou une durée minimale de maintien, et
compter les messages par côté dans les tests.
\iqlookfor{perte de priorité et hystérésis.}
\end{solution}

\begin{solution}{iq:hf:adverse-selection-and-toxicity:5}
Son obligation de cotation~: un temps minimal au meilleur prix ou à proximité et des largeurs maximales (livre~1, chapitre~24). On peut élargir dans la
largeur autorisée ou réduire la taille, mais ne pas s'effacer pendant longtemps.
\iqlookfor{les obligations comme contrainte sur la réponse.}
\end{solution}

\begin{solution}{iq:hf:adverse-selection-and-toxicity:6}
Exécutions restantes~: informées $\pi(1-a)$, non informées $(1-\pi)(1-b)$~; mark-out moyen $\bigl[(1-\pi)(1-b)G-\pi(1-a)L\bigr]/\bigl[\pi(1-a)+(1-\pi)(1-b)\bigr]$,
qui dépasse $(1-\pi)G-\pi L$ exactement lorsque $a>b$. Le volume baisse à $\pi(1-a)+(1-\pi)(1-b)$ de son niveau initial.
\iqlookfor{l'algèbre du mélange et la condition $a>b$.}
\end{solution}
